\documentclass[../../main.tex]{subfiles}
\graphicspath{{../../image/}{../../figure/}}
\begin{document}

\section{热弹性}\label{section:ch02-热弹性}

在连续介质热力学的框架中，一个\emph{热弹性介质}由如下本构假设来标识：对任一固定质点 $x$ 和任一运动，在 $x$ 和时刻 $t$ 处的内能 $\varepsilon$、Piola--Kirchhoff 应力 $S$、温度 $\theta$ 和参考热通量向量 $Q$ 之值，仅由 $(x,t)$ 处的变形梯度 $F$、熵 $s$ 和温度梯度 $G$ 之值通过本构方程确定：
\begin{align}
\begin{cases}
\varepsilon=\hat{\varepsilon}(F,s,G),\\
S=\hat{S}(F,s,G),\\
\theta=\hat{\theta}(F,s,G),\\
Q=\hat{Q}(F,s,G),
\end{cases}
\label{align:ch02-2-5-1}
\end{align}
其中 $\hat{\varepsilon},\hat{S},\hat{\theta},\hat{Q}$ 是定义在 $\{F\in\mathbb{M}^{m\times m}:\det F>0\}\times\mathbb{R}\times\mathbb{R}^{m}$ 子集上的光滑函数。此外 $\hat{\theta}(F,s,G)>0$。当热弹性介质\emph{均匀}时，同一组函数 $\hat{\varepsilon},\hat{S},\hat{\theta},\hat{Q}$ 以及同一参考密度值 $\rho_{0}$ 适用于所有质点 $x\in\mathscr{B}$。

Cauchy 应力 $T$ 和空间热通量 $q$ 也由相同形式的本构方程确定，它们可由\eqref{align:ch02-2-5-1} 结合\eqref{align:ch02-2-3-6}、\eqref{align:ch02-2-3-12} 导出。当采用运动的空间描述时，自然在清单\eqref{align:ch02-2-5-1} 中用 $T$ 和 $q$ 的本构方程替换 $S$ 和 $Q$ 的本构方程；同时在状态变量清单 $(F,s,G)$ 中用空间温度梯度 $g$（参见\eqref{align:ch02-2-3-17}）替换参考温度梯度 $G$。

上述本构方程必须满足前面所规定的条件。首先，如\cref{section:ch02-连续介质热力学平衡律} 所公设，任何平衡了质量、动量和能量的光滑热力学过程都必须恒等地满足 Clausius--Duhem 不等式\eqref{align:ch02-2-3-13}，或等价地，耗散不等式\eqref{align:ch02-2-3-15}。将\eqref{align:ch02-2-5-1} 代入\eqref{align:ch02-2-3-15}，得
\begin{align}
\mathrm{tr}\,[(\rho_{0}\partial_{F}\hat{\varepsilon}-\hat{S})\dot{F}^{\top}]+\rho_{0}(\partial_{s}\hat{\varepsilon}-\hat{\theta})\dot{s}+\rho_{0}\partial_{G}\hat{\varepsilon}\dot{G}-\hat{\theta}^{-1}\hat{Q}\cdot G\leqslant 0.
\label{align:ch02-2-5-2}
\end{align}

显然，通过适当控制体力 $b$ 和热量供给 $r$，可以构造出平衡了质量、动量和能量的光滑过程，并在某点 $(x,t)$ 处对 $F,s,G,\dot{F},\dot{s},\dot{G}$ 取任意预先指定的值，唯一约束是 $\det F>0$。因此\eqref{align:ch02-2-5-2} 不可能恒成立，除非本构关系\eqref{align:ch02-2-5-1} 取如下特殊形式：
\begin{align}
\begin{cases}
\varepsilon=\hat{\varepsilon}(F,s),\\
S=\rho_{0}\partial_{F}\hat{\varepsilon}(F,s),\\
\theta=\partial_{s}\hat{\varepsilon}(F,s),\\
Q=\hat{Q}(F,s,G),
\end{cases}
\label{align:ch02-2-5-3}
\end{align}
\begin{align}
\hat{Q}(F,s,G)\cdot G\geqslant 0.
\label{align:ch02-2-5-4}
\end{align}

于是内能可依赖于变形梯度和熵，但不依赖于温度梯度。应力和温度的本构方程由内能的本构方程通过\emph{caloric 关系}诱导，同样与温度梯度无关。只有热通量可依赖于温度梯度，但须满足条件\eqref{align:ch02-2-5-4}，它蕴涵热总是从物体较热部分流向较冷部分。

本构关系的另一项要求是遵守\cref{section:ch02-材料标架无差异} 中表述的材料标架无差异原理。将\eqref{align:ch02-2-4-4} 与 \eqref{align:ch02-2-4-3}$_{1}$ 结合\eqref{align:ch02-2-5-3}，我们推出函数 $\hat{\varepsilon}$ 和 $\hat{Q}$ 必须满足条件
\begin{align}
\hat{\varepsilon}(OF,s)=\hat{\varepsilon}(F,s),\qquad\hat{Q}(OF,s,G)=\hat{Q}(F,s,G),
\label{align:ch02-2-5-5}
\end{align}
对所有正常正交矩阵 $O$ 成立。一个简单计算可验证：当\eqref{align:ch02-2-5-5} 成立时，\eqref{align:ch02-2-4-4} 中其余条件将由\eqref{align:ch02-2-5-3}$_{2}$ 和 \eqref{align:ch02-2-5-3}$_{3}$ 自动满足。

为看出\eqref{align:ch02-2-5-5} 的含义，取 $O=R^{\top}$（其中 $R$ 是\eqref{align:ch02-2-1-6} 中的旋转张量），推出
\begin{align}
\hat{\varepsilon}(F,s)=\hat{\varepsilon}(U,s),\qquad\hat{Q}(F,s,G)=\hat{Q}(U,s,G).
\label{align:ch02-2-5-6}
\end{align}
显然反之亦然：若\eqref{align:ch02-2-5-6} 成立，则对任一正常正交矩阵 $O$，\eqref{align:ch02-2-5-5} 成立。因此，材料标架无差异原理完全编码在如下陈述中：内能和参考热通量向量可仅通过右拉伸张量 $U$ 依赖于变形梯度 $F$。

当采用运动的空间描述时，Cauchy 应力的本构方程
\begin{align}
T=\rho\,\partial_{F}\hat{\varepsilon}(F,s)F^{\top},
\label{align:ch02-2-5-7}
\end{align}
由\eqref{align:ch02-2-3-6}、\eqref{align:ch02-2-3-3} 和 \eqref{align:ch02-2-5-3}$_{2}$ 推出，只要\eqref{align:ch02-2-5-6} 成立，它就满足材料标架无差异原理\eqref{align:ch02-2-4-6}$_{1}$。对于空间热通量向量的本构方程
\begin{align}
q=\hat{q}(F,s,g),
\label{align:ch02-2-5-8}
\end{align}
材料标架无差异原理要求（回顾\eqref{align:ch02-2-4-6}$_{2}$、\eqref{align:ch02-2-4-3}$_{1}$ 和 \eqref{align:ch02-2-4-5}$_{2}$）：
\begin{align}
\hat{q}(OF,s,Og)=O\hat{q}(F,s,g),
\label{align:ch02-2-5-9}
\end{align}
对所有正常正交矩阵 $O$ 成立。

本构关系的最后一项一般要求是 Piola--Kirchhoff 应力满足角动量平衡\eqref{align:ch02-2-3-8}。但这并不施加额外限制，因为一个简单计算表明：一旦\eqref{align:ch02-2-5-5}$_{1}$ 成立，通过\eqref{align:ch02-2-5-3}$_{2}$ 算出的 $S$ 将自动满足\eqref{align:ch02-2-3-8}。因此在热弹性中，材料标架无差异蕴涵角动量平衡。

当介质被赋予\emph{材料对称性}时，本构方程进一步约化。回顾\cref{section:ch02-连续介质热力学平衡律}：当物体的参考构型通过一个等容双 Lipschitz 同胚 $\bar{x}$（其幺模 Jacobi 矩阵为 $H$，见\eqref{align:ch02-2-2-5}）改变时，场按照规则\eqref{align:ch02-2-3-19}、\eqref{align:ch02-2-3-20}、\eqref{align:ch02-2-3-21} 和 \eqref{align:ch02-2-3-22} 变换。特别地，由此可知，相对于原参考构型是热弹性的任何介质，相对于新参考构型也仍是热弹性的，尽管本构函数一般会改变。使 $\varepsilon,T,\theta$ 的本构函数保持不变的参考构型等容变换，体现了介质的\emph{材料对称性}。考察任一这样的变换，设其 Jacobi 矩阵为 $H$。由\eqref{align:ch02-2-3-19}$_{1}$、\eqref{align:ch02-2-3-20}$_{2}$ 和 \eqref{align:ch02-2-5-3}$_{1}$，内能的本构函数 $\hat{\varepsilon}$ 将保持不变，只要
\begin{align}
\hat{\varepsilon}(FH^{-1},s)=\hat{\varepsilon}(F,s).
\label{align:ch02-2-5-10}
\end{align}
一个简单计算可验证：当\eqref{align:ch02-2-5-10} 成立时，通过\eqref{align:ch02-2-5-7} 和 \eqref{align:ch02-2-5-3} 确定的 $T$ 和 $\theta$ 的本构函数自动在该 $H$ 下不变。由\eqref{align:ch02-2-3-19}$_{1}$ 和 \eqref{align:ch02-2-3-22}$_{2}$，热通量的本构函数 $\hat{q}$ 在 $H$ 下不变当且仅当
\begin{align}
\hat{q}(FH^{-1},s,g)=\hat{q}(F,s,g).
\label{align:ch02-2-5-11}
\end{align}

显然，满足\eqref{align:ch02-2-5-10} 和 \eqref{align:ch02-2-5-11} 的行列式为一的矩阵 $H$ 全体构成特殊线性群 $\mathrm{SL}(m)$ 的一个子群 $\mathscr{G}$，称为介质的\emph{对称群}。在某些介质中，$\mathscr{G}$ 可能只含单位矩阵 $I$，此时材料对称性是极小的。当 $\mathscr{G}$ 非平凡时，它通过\eqref{align:ch02-2-5-10} 和 \eqref{align:ch02-2-5-11} 对介质的本构函数施加条件。

当 $\mathscr{G}\equiv\mathrm{SL}(m)$ 时达到极大材料对称性。此时介质是一个\emph{热弹性流体}。取 $H=(\det F)^{-1/m}F\in\mathrm{SL}(m)$，应用\eqref{align:ch02-2-5-10} 和 \eqref{align:ch02-2-5-11}，我们推出 $\hat{\varepsilon}$ 和 $\hat{q}$ 可仅通过其行列式依赖于 $F$，或等价地，由\eqref{align:ch02-2-3-3}，通过密度 $\rho$ 依赖于 $F$：
\begin{align}
\varepsilon=\tilde{\varepsilon}(\rho,s),\qquad q=\tilde{q}(\rho,s,g).
\label{align:ch02-2-5-12}
\end{align}
Cauchy 应力可由\eqref{align:ch02-2-5-7} 得到，温度由\eqref{align:ch02-2-5-3}$_{3}$ 得到。计算给出
\begin{align}
T=-pI,
\label{align:ch02-2-5-13}
\end{align}
\begin{align}
p=\rho^{2}\partial_{\rho}\tilde{\varepsilon}(\rho,s),\qquad\theta=\partial_{s}\tilde{\varepsilon}(\rho,s).
\label{align:ch02-2-5-14}
\end{align}
在热力学标准教材中，\eqref{align:ch02-2-5-14} 通常以\emph{Gibbs 关系}的形式给出：
\begin{align}
\theta\,\mathrm{d}s=\mathrm{d}\varepsilon+p\,\mathrm{d}\!\left(\frac{1}{\rho}\right).
\label{align:ch02-2-5-15}
\end{align}

\eqref{align:ch02-2-5-12} 中的本构函数 $\tilde{q}$ 还必须满足材料标架无差异要求\eqref{align:ch02-2-5-9}，它现在取简单形式
\begin{align}
\tilde{q}(\rho,s,Og)=O\tilde{q}(\rho,s,g),
\label{align:ch02-2-5-16}
\end{align}
对所有正常正交矩阵 $O$ 成立。满足\eqref{align:ch02-2-5-16} 的 $\tilde{q}$ 的最终约化为
\begin{align}
q=\kappa(\rho,s,|g|)g,
\label{align:ch02-2-5-17}
\end{align}
其中 $\kappa$ 是一个标量值函数。于是我们已经证明：在热弹性流体中，内能仅依赖于密度和熵。Cauchy 应力是一个\emph{静水压强}，同样仅依赖于密度和熵。热通量服从 Fourier 定律，其\emph{热导率} $\kappa$ 可随密度、熵和热通量的大小而变化。

热弹性流体最简单的经典例子是\emph{理想气体}，它由 Boyle 定律
\begin{align}
p=R\rho\theta,
\label{align:ch02-2-5-18}
\end{align}
结合内能正比于温度的本构假设来标识：
\begin{align}
\varepsilon=c\theta.
\label{align:ch02-2-5-19}
\end{align}
在\eqref{align:ch02-2-5-18} 中，$R$ 是通用气体常数除以气体分子量；\eqref{align:ch02-2-5-19} 中的 $c$ 是比热。常数 $\gamma=1+R/c$ 是\emph{绝热指数}。经典分子运动论预言 $\gamma=1+2/n$，其中 $n$ 是气体分子的自由度数。当气体为单原子时达到最大值 $\gamma=5/3$。

将\eqref{align:ch02-2-5-18} 和 \eqref{align:ch02-2-5-19} 与 \eqref{align:ch02-2-5-13}、\eqref{align:ch02-2-5-14} 相结合，容易推出归一化单位下理想气体的本构关系为
\begin{align}
\varepsilon=c\rho^{\gamma-1}e^{s/c},\qquad p=R\rho^{\gamma}e^{s/c},\qquad\theta=\rho^{\gamma-1}e^{s/c}.
\label{align:ch02-2-5-20}
\end{align}

理想气体模型在很宽的密度和温度范围内令人满意地描述了普通气体的行为，但在状态变量取极端值时变得不太可靠，尤其在接近向液相转变的点附近。因此，人们提出了大量具有理论或经验来源的方程以适用于“实际气体”。最经典的例子是\emph{van der Waals 气体}，其中\eqref{align:ch02-2-5-18} 被替换为
\begin{align}
(p+a\rho^{2})(1-b\rho)=R\rho\theta,
\label{align:ch02-2-5-21}
\end{align}
其中 $a,b$ 是正参数。它对应本构关系
\begin{align}
\varepsilon=c\left(\frac{\rho}{1-b\rho}\right)^{\gamma-1}e^{s/c}+a\rho,\quad
p=R\left(\frac{\rho}{1-b\rho}\right)^{\gamma}e^{s/c},\quad
\theta=\left(\frac{\rho}{1-b\rho}\right)^{\gamma-1}e^{s/c}.
\label{align:ch02-2-5-22}
\end{align}

一个更奇特的模型是\emph{Chaplygin 气体}，其状态方程形如
\begin{align}
\varepsilon=\frac{1}{2\rho^{2}}f(s)-\frac{1}{\rho}g(s)+h(s),\quad
p=g(s)-\frac{1}{\rho}f(s),\quad
\theta=\frac{1}{2\rho^{2}}f'(s)-\frac{1}{\rho}g'(s)+h'(s).
\label{align:ch02-2-5-23}
\end{align}
注意在低密度下压强变为负，这与传统观念相悖。然而正是这一特征使 Chaplygin 气体对宇宙学家具有吸引力，因为他们将其与“暗物质”联系起来。

一个\emph{各向同性热弹性固体}是指对称群 $\mathscr{G}$ 为正常正交群 $\mathrm{SO}(m)$ 的热弹性材料。在该情形下，为得到内能函数 $\hat{\varepsilon}$ 的约化形式，我们将\eqref{align:ch02-2-5-10} 与 \eqref{align:ch02-2-5-6}$_{1}$ 相结合。回顾\eqref{align:ch02-2-1-7}，推出
\begin{align}
\hat{\varepsilon}(OUO^{\top},s)=\hat{\varepsilon}(U,s),
\label{align:ch02-2-5-24}
\end{align}
对任一正常正交矩阵 $O$ 成立。特别地，对使对称矩阵 $U$ 对角化的正常正交矩阵 $O$ 应用\eqref{align:ch02-2-5-24}：$OUO^{\top}=\Lambda$。这就确立了：由材料标架无差异和材料对称性，各向同性热弹性固体的内能可仅作为右拉伸张量 $U$ 特征值的对称函数依赖于 $F$。等价地，
\begin{align}
\varepsilon=\tilde{\varepsilon}(J_{1},\cdots,J_{m},s),
\label{align:ch02-2-5-25}
\end{align}
其中 $(J_{1},\cdots,J_{m})$ 是 $U$ 的不变量。特别地，当 $m=3$ 时，可取 $J_{1}=|F|^{2}$，$J_{2}=|F^{*}|^{2}$，$J_{3}=\det F$，其中 $F^{*}$ 是 $F$ 的伴随矩阵。各向同性热弹性固体的 Cauchy 应力约化形式由\eqref{align:ch02-2-5-25} 和 \eqref{align:ch02-2-5-7} 计算，记录在\cref{section:ch02-注记} 所引用的文献中。读者还可在那里找到具体可压缩或不可压缩各向同性弹性固体本构函数的明确例子。

在热弹性的另一种等价表述中，人们将温度 $\theta$ 而非熵 $s$ 作为状态变量，并写出关于 $s$ 而非关于 $\theta$ 的本构方程。此时监测\emph{Helmholtz 自由能}
\begin{align}
\psi=\varepsilon-\theta s
\label{align:ch02-2-5-26}
\end{align}
来代替内能 $\varepsilon$ 也是方便的。于是人们从如下本构方程出发：
\begin{align}
\begin{cases}
\psi=\bar{\psi}(F,\theta,G),\\
S=\bar{S}(F,\theta,G),\\
s=\bar{s}(F,\theta,G),\\
Q=\bar{Q}(F,\theta,G),
\end{cases}
\label{align:ch02-2-5-27}
\end{align}
以代替\eqref{align:ch02-2-5-1}。要求所有平衡了质量、动量和能量的光滑热力学过程都必须恒等地满足耗散不等式\eqref{align:ch02-2-3-15}，这一要求将\eqref{align:ch02-2-5-27} 约化为
\begin{align}
\begin{cases}
\psi=\bar{\psi}(F,\theta),\\
S=\rho_{0}\partial_{F}\bar{\psi}(F,\theta),\\
s=-\partial_{\theta}\bar{\psi}(F,\theta),\\
Q=\bar{Q}(F,\theta,G),
\end{cases}
\label{align:ch02-2-5-28}
\end{align}
\begin{align}
\bar{Q}(F,\theta,G)\cdot G\geqslant 0,
\label{align:ch02-2-5-29}
\end{align}
它们是\eqref{align:ch02-2-5-3}、\eqref{align:ch02-2-5-4} 的类比\footnote{形如\eqref{align:ch02-2-5-3} 的本构方程称为 caloric 形式，形如\eqref{align:ch02-2-5-28} 的称为 thermal 形式。}。材料标架无差异原理和材料对称性的存在进一步约化上述本构方程。特别地，$\bar{\psi}$ 满足与上面 $\hat{\varepsilon}$ 相同的条件。

我们以关于特殊热力学过程的注记来结束热弹性的讨论。若热通量 $Q$ 恒等于零，则过程称为\emph{绝热的}；若温度场 $\theta$ 为常数，则过程称为\emph{等温的}；若熵场 $s$ 为常数，则过程称为\emph{等熵的}。注意\eqref{align:ch02-2-5-29} 蕴涵 $\bar{Q}(F,\theta,0)=0$，因此特别地，所有等温过程都是绝热的。导热性差的材料通常被建模为热的\emph{不良导体}，其特征是本构假设 $\hat{Q}\equiv 0$。于是不良导体的每一热力学过程都是绝热的。

在等熵过程中，熵取为常数 $s\equiv\bar{s}$；温度和热通量的本构关系被弃去，内能和应力的本构关系限制在 $s=\bar{s}$ 上：
\begin{align}
\begin{cases}
\varepsilon=\hat{\varepsilon}(F,\bar{s}),\\
S=\rho_{0}\partial_{F}\hat{\varepsilon}(F,\bar{s}).
\end{cases}
\label{align:ch02-2-5-30}
\end{align}
特别地，对理想气体，由\eqref{align:ch02-2-5-20}，
\begin{align}
\varepsilon=\frac{\kappa}{\gamma-1}\rho^{\gamma-1},\qquad p=\kappa\rho^{\gamma},
\label{align:ch02-2-5-31}
\end{align}
其中 $\kappa=R\exp(\bar{s}/c)$。

在等熵过程中，运动仅由质量和动量平衡律结合本构关系\eqref{align:ch02-2-5-30} 确定。这可能给人一种印象：等熵热弹性同构于纯力学的\emph{超弹性}理论。然而这不完全准确，因为等熵热弹性从热力学继承了第二定律，其形式为：要维持一个等熵过程，必须以如下方式控制热量供给 $r$，使得在常数熵场下随之而来的运动满足能量平衡律\eqref{align:ch02-2-3-10}。当过程同时也是绝热的（$Q=0$）时，Clausius--Duhem 不等式\eqref{align:ch02-2-3-13} 化为 $r\leqslant 0$，此时\eqref{align:ch02-2-3-10} 蕴涵
\begin{align}
(\rho_{0}\varepsilon+\tfrac{1}{2}\rho_{0}|v|^{2})^{\cdot}\leqslant\mathrm{Div}\,(v^{\top}S)+\rho_{0}v^{\top}b.
\label{align:ch02-2-5-32}
\end{align}
该不等式的 Euler 形式为
\begin{align}
(\rho\varepsilon+\tfrac{1}{2}\rho|v|^{2})_{t}+\mathrm{div}\,[(\rho\varepsilon+\tfrac{1}{2}\rho|v|^{2})v^{\top}]\leqslant\mathrm{div}\,(v^{\top}T)+\rho v^{\top}b.
\label{align:ch02-2-5-33}
\end{align}

上述不等式在等熵热弹性中所起的作用，正如 Clausius--Duhem 不等式\eqref{align:ch02-2-3-13}、\eqref{align:ch02-2-3-14} 在一般热弹性中所起的作用：对光滑运动，它们由\eqref{align:ch02-2-3-4} 和 \eqref{align:ch02-2-5-30} 以等式恒成立；相反，在运动仅仅 Lipschitz 连续的情形中，它们是作为运动\emph{热力学可容许性}检验的额外条件。

实际上，当判断熵涨落的影响可忽略时，就采用等熵理论。这并非罕见情形，理由如下：在光滑绝热过程中 $Q=0$，若此外还有 $r=0$，则\eqref{align:ch02-2-3-18} 结合\eqref{align:ch02-2-5-3} 给出 $\dot{s}=0$。于是在无热量供给时，从等熵出发的绝热过程只要是光滑的就保持等熵。当 $F,v,s$ 仅仅 Lipschitz 连续时，光滑性要求即告满足，这允许带弱间断面的过程，但不允许激波。正如我们后面将看到的，即使激波已经形成，只要它们保持弱，与密度和速度的涨落相比，熵涨落是小量（三阶），因而可忽略。

在等温热弹性中，$\theta$ 取为常数 $\hat{\theta}$，热量供给 $r$ 被调节以平衡能量方程，运动仅由质量和动量平衡律确定。所需的唯一本构方程是
\begin{align}
\begin{cases}
\psi=\bar{\psi}(F,\hat{\theta}),\\
S=\rho_{0}\partial_{F}\bar{\psi}(F,\hat{\theta}),
\end{cases}
\label{align:ch02-2-5-34}
\end{align}
即\eqref{align:ch02-2-5-30} 的类比。热力学第二定律的推论如前一样，通过将\eqref{align:ch02-2-3-10} 与 \eqref{align:ch02-2-3-13} 相结合来考察，此时假定 $\theta=\hat{\theta}=$ 常数。这给出
\begin{align}
(\rho_{0}\psi+\tfrac{1}{2}\rho_{0}|v|^{2})^{\cdot}\leqslant\mathrm{Div}\,(v^{\top}S)+\rho_{0}v^{\top}b,
\label{align:ch02-2-5-35}
\end{align}
其 Euler 形式为
\begin{align}
(\rho\psi+\tfrac{1}{2}\rho|v|^{2})_{t}+\mathrm{div}\,[(\rho\psi+\tfrac{1}{2}\rho|v|^{2})v^{\top}]\leqslant\mathrm{div}\,(v^{\top}T)+\rho v^{\top}b,
\label{align:ch02-2-5-36}
\end{align}
应与\eqref{align:ch02-2-5-32} 和 \eqref{align:ch02-2-5-33} 相比较。我们得出结论：等温和等熵热弹性本质上是同构的，前者中恒温下的 Helmholtz 自由能扮演后者中恒熵下内能的角色。

对等温过程，理想气体 $\theta=\hat{\theta}$，
\begin{align}
\psi=k\log\rho,\qquad p=k\rho,
\label{align:ch02-2-5-37}
\end{align}
其中 $k=R\hat{\theta}$。

\end{document}
